\subsection{Auswahl des Projektprofils}
\label{subsec:auswahl-projektprofil}

Die Zugriffskanäle bestimmen das Profil: \texttt{web-only} steht für Browser
ohne Backend, \texttt{web-cloud} für Browser mit Backend,
\texttt{desktop-local} für Desktop ohne Cloud-Backend,
\texttt{desktop-cloud} für Desktop mit Backend und \texttt{full-platform} für
Browser und Desktop gemeinsam. Ein Profil für Backend ohne providerneutrale
Deployment-Dateien existiert nicht.

Gewählt wird das kleinste passende Profil. Erst danach wird PostgreSQL als
separate Achse geprüft und nur bei Bedarf zu einem Backendprofil ergänzt. Für
\texttt{web-only} und \texttt{desktop-local} ist dies unzulässig
\parencite{kleiveist2026profiles}.

Dokumentiert wird, ob das Profil vollständig passt, benannte Erweiterungen
benötigt oder zentrale Anforderungen verfehlt. Im letzten Fall sind eine
andere Architektur oder klar abgegrenzte Teilnutzung zu prüfen, nicht
ersatzweise das größte Profil.
Teamkompetenz und wirtschaftliche Randbedingungen ergänzen diese technische
Entscheidung, ändern aber nicht die Bedeutung der Profile.
\beginnerinput{workspace/statement/800_einsatz_transfer/840}
